\chapter{Considerações finais}\label{cap:consideracoes-finais}

Este capítulo encerra o relatório apresentando a síntese das atividades desenvolvidas durante o estágio, as dificuldades encontradas no percurso, as limitações identificadas no trabalho realizado e as perspectivas de continuidade do projeto.

\section{Síntese das atividades desenvolvidas}

O estágio teve como resultado o desenvolvimento do S.R.I. -- Sistema de Rotas Inteligentes, aplicativo móvel voltado ao planejamento, à execução e à reutilização de rotas de coleta e entrega. O trabalho foi conduzido segundo o processo RUP, o que implicou a produção dos artefatos de visão, de levantamento de requisitos, de especificação de caso de uso, de especificação complementar e de glossário, além dos modelos de casos de uso, de classes, de entidade e relacionamento, de sequência, de estados e de implantação \cite{kruchten2003}.

Do ponto de vista da implementação, a adoção do framework Flutter permitiu manter uma única base de código para as plataformas Android e iOS, enquanto a persistência local em SQLite assegurou o funcionamento das operações de cadastro e de consulta independentemente da disponibilidade da conexão de rede \cite{flutter2024, sqlite2024}. A integração com os serviços do OpenStreetMap viabilizou a geocodificação de endereços e a representação cartográfica dos pontos cadastrados sem a contratação de chaves de acesso comerciais \cite{osm2024}.

O caso de uso de gerenciamento de locais foi priorizado como primeiro escopo de entrega, por constituir pré-requisito para as funcionalidades de geração de rotas e de rotas salvas \cite{sriEspecUC2026}. A opção por essa ordem de implementação permitiu validar, desde o início do desenvolvimento, a integração entre a interface, a camada de persistência e o serviço externo de geocodificação, que é a base tecnológica compartilhada por todas as demais funcionalidades do sistema. A partir dessa base, foram implementados, na sequência, os casos de uso de geração de rotas, de gerenciamento de rotas salvas e de gerenciamento de usuário, além da navegação assistida, de modo que a versão consolidada contempla a totalidade dos casos de uso previstos no documento de especificação.

\section{Resultados alcançados}

A versão consolidada do S.R.I. contempla os quatro casos de uso previstos na especificação e a funcionalidade de navegação assistida. O gerenciamento de locais permite criar, consultar, atualizar e excluir pontos de interesse, com conversão automática do endereço em coordenadas. A geração de rotas permite compor trajetos nas modalidades ponto a ponto e em circuito, com ordenação automática das paradas, apresentação do traçado sobre o mapa e indicação da distância total e da duração estimada. O gerenciamento de rotas salvas permite gravar, reabrir, alterar e excluir os trajetos calculados. O gerenciamento de usuário impõe a autenticação prévia exigida pelos requisitos e responde pela criação da conta, pela manutenção dos dados cadastrais, pela troca de senha e pela exclusão da conta.

A verificação automatizada produziu resultados objetivos. A análise estática do código-fonte, executada com as regras de estilo configuradas no projeto, não registrou nenhum erro nem aviso. A suíte de testes automatizados, composta por noventa e cinco testes, foi executada integralmente com aprovação. A compilação completa do aplicativo foi concluída com sucesso para o alvo web, confirmando que o conjunto do código-fonte compila de ponta a ponta. O código do aplicativo totaliza 6.047 linhas distribuídas em 37 arquivos, acrescidas de 1.652 linhas de testes.

Quanto à qualidade do algoritmo de ordenação das paradas, o teste comparativo com a solução ótima obtida por força bruta, em uma instância de cinco pontos, confirmou que a solução heurística não excede o ótimo em mais de quinze por cento, além de nunca ser pior que a ordem produzida isoladamente pela heurística construtiva.

Duas métricas de aceitação declaradas nos requisitos permanecem pendentes de medição em campo. A primeira é o tempo de resposta das operações de geocodificação e de cálculo de trajeto, que o projeto limita em trinta segundos: o limite de tempo das requisições foi configurado em vinte segundos justamente para assegurar esse teto, mas a medição efetiva depende da latência real dos serviços externos em condições de uso. A segunda é o patamar de assertividade de setenta por cento nas rotas geradas, que depende da validação das rotas junto ao usuário do sistema e não pode ser aferido apenas por testes automatizados \cite{sriEspecComp2025, sriPedido2025}.

\section{Dificuldades encontradas}

A primeira dificuldade decorreu da assimetria entre o esforço de documentação e o tempo disponível para implementação. A adoção do RUP, com a produção de artefatos formais em cada disciplina, mostrou-se valiosa para o entendimento do domínio, mas exigiu um volume de trabalho de escrita que competiu diretamente com as horas destinadas ao desenvolvimento. O aprendizado obtido foi o de dimensionar o nível de detalhe da documentação conforme o risco do requisito, aprofundando os artefatos das funcionalidades centrais e mantendo mais enxutos aqueles de menor incerteza \cite{kruchten2003}.

A segunda dificuldade esteve relacionada à integração com os serviços externos de geocodificação e de roteamento. Tais serviços impõem limites de uso e apresentam variação no tempo de resposta, o que exigiu o tratamento explícito de situações de falha e de espera prolongada na interface, com a exibição de indicadores de carregamento e de mensagens de erro compreensíveis ao usuário \cite{nominatim2024, osrm2024}.

A terceira dificuldade foi de natureza algorítmica. A definição da estratégia de ordenação das paradas exigiu o estudo das heurísticas de roteirização e a avaliação de seu custo computacional no dispositivo, uma vez que a obtenção da solução ótima é inviável para instâncias de tamanho realista \cite{cunha2000}. A solução adotada combinou uma heurística construtiva com um método de refinamento, o que demandou a construção de testes específicos para verificar a corretude da ordenação produzida \cite{croes1958}.

\section{Limitações}

A principal limitação identificada é a dependência de conectividade para as funcionalidades que exigem geocodificação de endereços e cálculo de rota sobre a malha viária. O cadastro e a consulta de locais permanecem operacionais sem conexão, mas a conversão de um novo endereço em coordenadas e a obtenção de um trajeto por vias exigem acesso à rede, conforme estabelecido nas premissas e dependências do produto \cite{sriVisao2025}.

Uma segunda limitação decorre do escopo funcional da navegação assistida. A funcionalidade mantém o motorista informado sobre a próxima parada e o quanto falta para alcançá-la, mas não recalcula o trajeto quando o veículo se desvia da malha viária prevista nem emite comandos de voz, recursos presentes em soluções comerciais de navegação. A verificação funcional em aparelho físico, com variação real da precisão do sinal de GPS e instabilidade da rede móvel, também permanece pendente, tendo a validação automatizada se limitado aos cálculos de navegação e às regras de negócio \cite{sriEspecComp2025}.

Registra-se ainda que os testes foram executados em ambiente de emulador, o que, embora permita verificar o comportamento das telas, das permissões de localização e da persistência, não reproduz integralmente as condições de uso em campo, como a variação da precisão do sinal de GPS e a instabilidade da rede móvel.

\section{Trabalhos futuros}

Como continuidade do projeto, estão previstas a medição das métricas de aceitação ainda pendentes, com o registro dos tempos de resposta das operações de geocodificação e de cálculo de trajeto e com a validação das rotas geradas junto ao usuário do sistema, de modo a consolidar a avaliação de assertividade \cite{sriPedido2025}; a realização de testes em dispositivo físico, em condições reais de operação, para aferir o comportamento da geolocalização, o consumo de bateria e o desempenho com o banco de dados já populado; e a ampliação da suíte de testes automatizados, com a cobertura dos fluxos de interface hoje verificados apenas manualmente.

No curto prazo, uma evolução de baixo custo é o emprego da geometria já gravada no cálculo da distância restante durante a navegação, em lugar da distância em linha reta hoje utilizada, o que aproximaria a estimativa apresentada ao motorista da distância efetivamente percorrida sobre as vias.

No médio prazo, o documento de Visão prevê a inclusão de recursos de visualização do tráfego em tempo real, o compartilhamento de rotas por link e a integração com aplicativos de entrega e com sistemas de gestão logística \cite{sriVisao2025}. Outra evolução considerada é a substituição da heurística de ordenação por métodos de busca local mais elaborados, com o objetivo de reduzir a distância total dos circuitos gerados.

\section{Considerações finais}

O desenvolvimento do S.R.I. proporcionou a aplicação integrada dos conteúdos estudados ao longo do curso, articulando engenharia de software, modelagem de sistemas, banco de dados, algoritmos e programação para dispositivos móveis em torno de um problema real, apresentado por um usuário real e com critérios de aceitação declarados. A experiência evidenciou que a qualidade do resultado depende menos da escolha da tecnologia do que da clareza dos requisitos e da consistência entre o processo de negócio documentado e as funcionalidades efetivamente implementadas.

Do ponto de vista formativo, o estágio consolidou a compreensão de que o desenvolvimento de software é uma atividade de engenharia, na qual decisões de escopo, de modelagem e de validação precisam ser tomadas de forma explícita e registrada. A documentação produzida ao longo do período, além de atender às exigências acadêmicas do estágio obrigatório, constitui o principal instrumento de continuidade do projeto, uma vez que permite que o trabalho seja retomado sem a necessidade de reconstruir o entendimento do domínio \cite{sommerville2019}.
